\providecommand{\PushA}{\frac{5422587}{6103515625}}
\providecommand{\PushK}{890666}
\providecommand{\PushTheta}{0.99999999997327}
\section{The comparison theorem: standing data, prices, and hypotheses}
\label{sec:comparison-theorem}

\paragraph{What this section proves and how to read it.}
This section states Theorem~C: if distributions $\mu,\lambda,\nu$
satisfy explicit hypotheses (CP0)--(CP8) and the number $\Theta$
formed from their prices is $<1$, then
$s(N)\le C'N^{1-a/(72-\theta)}$ with $0<\theta\le1$. No particular
local law is used; Part~V later exhibits laws satisfying the
hypotheses. In words: $\mu$ is the reflection-positive reference
law; $\lambda$ is obtained from $\mu$ by record differences with the
same conditional one-half distributions, grouped by the cut at which
each difference has a uniform operator bound (a \emph{tag}); a
positive law $\Nbr$ dominates $\lambda$ at the reserved cut in the
$2\times2$ block sense; $\nu$ is a finite signed combination of
nonnegative laws each accepting the queries (in the bundled sense of
\cref{subsec:intro:glossary}), with symmetric and skew differences
priced separately; second positive laws dominate $\nu$ itself.
Inputs: Part~II only. Output: \cref{thm:comparison}, proved in
\cref{sec:comparison-lemmas,sec:comparison-proof} and applied with
$k=6$ in Part~V (\cref{tab:hypotheses}). A first reading should
follow (CP0), the price definitions, the statement of Theorem~C with
the reading guide preceding it, and treat (CP1)--(CP8) as a checklist
to be consulted in Part~V. Two reading-order warnings: the
repricing corollary in \cref{subsec:cp:prices} mentions totals
$L_0,B$ defined in (CP2)--(CP3) below and can be skipped until then;
and \emph{robust} bounds, used in (CP5) and (CP7), are defined in
\cref{def:robust} immediately after (CP7)---in brief, bounds whose
$o(1)$ errors come only from the two allowed analytic remainders.
Theorem~C is the $k$-query form of \cref{prop:mech:one-query}.

\emph{Letters.} In this part $B$ is the energy factor of the law
dominating $\lambda$, $A$ the total skew price, $\mathcal E$ the
dominated price bound for $\nu-\lambda$ and $J$ the coupled index
term; an
atom (a block of primes) is written $\Blk$. The square-difference-free set of
integers of the main theorem is also written $A$, but only in its lift $F_A$,
as the second entry of a pair $(\Lsc,A)$ (in $\Yf(\Lsc,A)$, in a
counterexample, and in ``uniformly in $A$''), and where such a set is
quantified (as in \cref{def:hyb:coupled} and Theorem~C); its density is
$\alpha=|A|/N$.

\subsection{Standing data (CP0), prices, and conversion between exponents}\label{subsec:cp:prices}

Theorem~C is stated for abstract comparison data, which Part~V
later verifies for one concrete construction. This subsection lists
those data as (CP0), then defines the price of changing the law at
one atom and how a price computed at a smaller exponent $a_0$ is
converted to the exponent $a$ of the theorem.

\begin{definition}[(CP0) Standing data: slots, scales, queries, index]\label{def:cp0}
\emph{Comparison data} consist of the following objects.
\begin{enumerate}[label=(\alph*)]
\item\emph{Slots and cuts.} The slot array $V$ of \cref{def:slot-array},
 $d=|V|=72$, with the family $\mathfrak F$ of its cuts: the chamber cuts
 $\vartheta_{c,a}$ of the two coordinates and the reserved cut
 $\iota(u_0,u_1,\ell)=(u_0,u_1,1-\ell)$, with plus half $V_+=\{\ell=0\}$ and
 minus half $V_-=\iota V_+$. The family contains the sorting sequence of
 \cref{lem:slot-array}(iii) (\cref{def:sorting}).
\item\emph{Scales, atoms and the analytic stage.} A fixed integer $\Kint\ge2$
 and, for every scale $\Lsc\ge1$ with $N=\Kint\Lsc$, the group $G_N$, its set
 $\mathcal A_N$ of atoms, the lift $F_A$ and the inputs $g^{I,D_0,s}$ of
 \cref{def:record:functional}, with fixed prefix exponents $e_p$; the
 analytic stage of \cref{def:record:analytic-stage}; a density exponent
 $0<\theta\le1$; and a price exponent $a$ with $0<a<\sigma$. We write
 $\mathcal A=\bigcup_N\mathcal A_N$ for the set of atoms of all scales. Every
 law attached below to an atom $\Blk$ is a measure on $X_{\Blk}^V$ that does
 not depend on the scale.
\item\emph{Queries.} A family $\mathcal E_k$ of $k\ge1$ directed queries
 $e=(u_e,v_e)$ with pairwise disjoint endpoints (the $2k$ slots $u_e,v_e$ are
 distinct), sources $u_e\in V_+$ and reflected targets $\iota v_e\in V_+$.
\item\emph{Index law.} A fixed modulus $\Mzero=\prod_p p^{h_p}$,
 $h_2\ge3$, coprime to every prime of $G_N$; a product convolution law
 on nonvanishing square increments at these coordinates; and independent
 conductor-depth envelopes and moment $R_I$, as in
 \cref{def:hyb:coupled}. Each exact-conductor space is preserved by the
 coordinate transition, and the norm of the transition restricted
 to it is bounded by the corresponding decreasing depth tail. The interval transition is the cyclic shift.
\end{enumerate}
\end{definition}

Only (a) refers to the particular array: the arguments of this part use of
$V$ only that its family of cuts contains a sorting sequence and the reserved
cut, and that the chamber folds commute with $\iota$ ((O4) of
\cref{def:admissible}; \cref{lem:slot-array}(iii)). So the scheme of
(CP1)--(CP8) is not tied to this array, and \cref{subsec:history:216} reads
the $216$-slot design, with $V=(\Zsix)^3$ and the reserved cut
$\vartheta_{2,0}$, as its first instance. The analytic stage of Part~II,
however, is formulated for $72$ slots, and Theorem~C is stated for the array
of (a), so $d=72$ throughout; the letter $d$ marks where the number of slots
enters an exponent (the slot price exponent $a/d$ and the density factor
$\theta/d$). In Part~V,
$k=6$ and $\mathcal E_6$ is the family of the six queries of
\cref{def:slot-array}(e), which satisfies (c) by \cref{lem:slot-array}(iv).

\paragraph{Energies, the counterexample and $o(1)$.}
For a finite signed measure $\kappa=\bigotimes_{\Blk}\kappa_{\Blk}$ on
$G_N^V$ let $C_\iota(\kappa)$ be its cut matrix at $\iota$ and
$Q_\kappa(X,Y)=\ip X{C_\iota(\kappa)Y}$ its cut form for half-functions $X,Y$
on $G_N^{V_+}$ (\cref{def:cc:order}), and put
\[
 Y_\kappa=Q_\kappa(\bar X,\bar X)=\Lambda_\kappa(F_A,\dots,F_A),
\]
where $\bar X$ is the half-function carrying $F_A$ at every slot of $V_+$.
Since every slot carries the same input, $Y_\kappa$ is also the value at
$(\bar X,\bar X)$ of the cut form of $\kappa$ at any other cut. For the
reference $\mu$ of (CP1), $Y_\mu=\Yf(\Lsc,A)$ and $\alpha^d\le Y_\mu$
(\cref{lem:record:holder}). Every argument of this part takes place at a
least normalized counterexample $(\Lsc,A)$ for $(a,\theta,C_0)$
(\cref{def:record:counterexample}), and $Y=Y_\mu>0$ there. Energies whose
selected atoms have a union $E$ with $q_E\le N^\tau$ are structured energies,
bounded by their coarse prices (\cref{prop:record:density-structured}), with
sharp bounds applied in the order of \cref{lem:record:folded-energy}. Every
expansion over selected atoms also has terms with $q_E>N^\tau$; their sum is
bounded by
\cref{lem:record:density-lp}, with the post-reduction prices of the selected
laws as the coefficients $\eta_{\Blk}$, and is $O(\Lsc^{-\sigma}\alpha)=o(1)Y$. Throughout,
$o(1)$ denotes a quantity that tends to zero as $\Lsc\to\infty$, uniformly in
$A$. Every $o(1)$ of this part is produced by structured energies and by the
two kinds of analytic remainders, the large selections just described and
the cancellations of \cref{lem:record:density-can}; this is the robustness of
\cref{def:robust}, which is what the hypothesis (Rec) of
\cref{thm:record:lineardensity} asks for (\cref{lem:record:rec}).

At a least counterexample, fixing the labels at an atom $\Blk$ fixes residue
classes, and splitting a class into progressions whose step is a square
divisible by $q_{\Blk}$ returns the problem at a shorter scale
(\cref{lem:record:fibre-scale}); by the induction hypothesis this costs the
factor $(q_{\Blk}^2)^{a/d}$ at each of the $d$ slots. This is the origin of
the prices.

\begin{definition}[Prices]\label{def:price}
Let $\Blk$ be an atom and $q_{\Blk}=\prod_{p\in\Blk}p$ its block modulus. Its
\emph{selection factor} is $q_{\Blk}^{2a}$; for a prefix atom $\{p\}$ with
$e_p$ even we read $q_{\Blk}^{2a}$ as $p^{e_pa}$, the factor of selecting the
full group $\Z/p^{e_p}\Z$ (\cref{def:record:structured-energy}).
\begin{enumerate}[label=(\alph*)]
\item The \emph{price} of a finite signed measure $\zeta$ on $X_{\Blk}^V$ is
 $\|\zeta\|_{\rm var}\,q_{\Blk}^{2a}$: selecting $\Blk$ changes every slot.
\item Let $c\in\mathfrak F$ be a cut, $\kappa$ a nonnegative measure on
 $X_{\Blk}^V$ of mass $w$ and $\zeta$ a signed measure on $X_{\Blk}^V$. Then
 $\zeta$ has \emph{sharp bound $d_\zeta$ relative to $\kappa$ at $c$} if, for a
 probability law $\pi$ on a finite set of shared labels,
 $\kappa=w\int\kappa_s\,\pi(ds)$ and $\zeta=w\int K_s\,\pi(ds)$ with
 probability measures $\kappa_s$ and signed measures $K_s$ that satisfy
 \eqref{eq:record:sharp} at $c$ relative to the half marginals of $\kappa_s$,
 with the constant $d_\zeta$, for every $s$. The \emph{sharp price} of $\zeta$
 is $w\,d_\zeta\,q_{\Blk}^{2a}$. More generally an \emph{operator price} is a
 weight, times an operator norm relative to actual half marginals, times
 $q_{\Blk}^{2a}$.
\item An \emph{upper price} is any number at least the price. A \emph{total}
 is a sum of prices over the atoms of $\mathcal A$. In an expansion over
 selected atoms each selected atom is charged its price once per expanded
 term.
\end{enumerate}
\end{definition}

The unpriced part of a price (a mass, a weight, a weighted norm) does not
depend on the exponent, and an atom's laws do not depend on the scale; so a
total over $\mathcal A$ bounds the corresponding sum over $\mathcal A_N$ at
every scale.

\begin{lemma}[Repricing]\label{lem:record:repricing}
Let $0<a'\le a''$ and $L\ge0$, and let $\mathcal P$ be a family of atoms of
primes with $e_p=1$ such that $\log q_{\Blk}\le L$ for every
$\Blk\in\mathcal P$. Then $q_{\Blk}^{2a''}\le e^{2(a''-a')L}q_{\Blk}^{2a'}$ for
$\Blk\in\mathcal P$, and for all numbers $c_{\Blk}\ge0$ that do not depend on
the exponent,
\[
 \sum_{\Blk\in\mathcal P}c_{\Blk}\,q_{\Blk}^{2a''}
 \le e^{2(a''-a')L}\sum_{\Blk\in\mathcal P}c_{\Blk}\,q_{\Blk}^{2a'} .
\]
In particular every price, sharp price or operator price of a law of an atom
of $\mathcal P$ at the exponent $a''$ is at most $e^{2(a''-a')L}$ times its
value at $a'$; an upper bound at $a'$ for a sum of such prices (for instance
the block charge \eqref{eq:record:cell-charge} of Part~V, in the notation
there), multiplied by $e^{2(a''-a')L}$, bounds the
sum at $a''$; and two such steps compose to one step from the first to the
last exponent.
\end{lemma}

\begin{proof}
For each atom, $q_{\Blk}^{2a''}=q_{\Blk}^{2a'}e^{2(a''-a')\log q_{\Blk}}\le
q_{\Blk}^{2a'}e^{2(a''-a')L}$, because $a''-a'\ge0$ and $\log q_{\Blk}\le L$.
Multiply by $c_{\Blk}\ge0$ and sum. A price is its unpriced part times
$q_{\Blk}^{2a}$, and the partition and the laws are the same at both
exponents; only the selection factor changes. Two steps with budgets $L_1$,
$L_2$ multiply to the factor $e^{2(a_1-a_0)L_1+2(a_2-a_1)L_2}$, which is
$e^{2(a_2-a_0)L}$ for a common budget $L$.
\end{proof}

The argument needs $a''\ge a'$: for a negative difference the upper bound
$L$ for $\log q_{\Blk}$ would point the wrong way.

\begin{corollary}[The first stage at a larger exponent]
\label{cor:first-stage-repricing}
Let $a_0\le a$, and let every first-stage atom of (CP2) be a block of primes
with $e_p=1$ and $\log q_{\Blk}\le L$. Let $U_0'$, $P'$ and $d_j'$, $v_j'$
($j\in\mathfrak T$) be upper bounds for the totals $U_0$, $P$ and the tag
totals $d_j$, $v_j$ of (CP2) and (CP3), computed for the same laws at the
exponent $a_0$, and put $t=e^{2(a-a_0)L}$. Then at the exponent $a$
\[
 U_0\le tU_0',\qquad P\le tP',\qquad d_j\le td_j',\qquad v_j\le tv_j'
 \quad(j\in\mathfrak T),
\]
and
\[
 L_0\le e^{tP'}\tagser(td',tv'),\qquad
 R_{\rm err}\le e^{tU_0'}\tagser(td',tv'),\qquad
 B\le e^{tU_0'}\bigl(1+\tagser(td',tv')\bigr).
\]
\end{corollary}

\begin{proof}
Every first-stage price is a price of a law of a first-stage atom: $p_\Blk$
and $u_\Blk$ are prices of positive records, $d_{\Blk,i}$ is a sharp price and
$v_{\Blk,i}$ the price of two positive records. So
\cref{lem:record:repricing}, applied to the sums over the first-stage atoms
and, for the tag totals, to the sums over the differences of each tag, gives
the first four bounds. The series $\tagser$ and the exponential are
nondecreasing in every argument (\cref{rem:cp:tags}), which gives the bounds
for $L_0=e^P\tagser(d,v)$ and $R_{\rm err}=e^{U_0}\tagser(d,v)$, and then
$B=e^{U_0}+R_{\rm err}\le e^{tU_0'}+e^{tU_0'}\tagser(td',tv')$.
\end{proof}

In the closing of Part~V the corollary is applied with the totals of the
middle range certified at the anchor $a_0=\PushAnchor$
(\cref{lem:record:middle-prices,cert:record:middle-price}) and the actual
middle block budget $L=570$ (\cref{def:closing-quantities}); the upper prices
have their own anchor $a_{\rm u}=883/10^6$ and are repriced band by band, each
with its own logarithmic product bound (\cref{prop:record:upper}).

\subsection{Queries and the random interval-and-residue index}\label{subsec:cp:index}

The primes dividing $\Mzero$ are not coordinates of $G_N$. They enter
through a fixed external residue index on each query, together with an
interval index.
The coupled index is $(b_e+1,s_e+r_e)$, where $r_e$ has the prescribed
product law on nonvanishing square classes. On a nonwrapping interval,
accepted-square cancellation follows from \cref{lem:record:density-can}
and the prime-power kernel extension \cref{prop:record:kernel20}(iii).
Uniform averaging still gives exactly the original lift. The residue depths
enter the estimates through the conductor moment $R_I$
(\cref{def:hyb:coupled}(g)).
The following elementary unit-square case also explains the earlier
squarefree index and remains useful when bounding its moment.

\begin{lemma}[Unit squares]\label{lem:func:unitsq}
Let $\Mzero=8d_0$ with $d_0$ odd and squarefree, and let $\Qsqmod{\Mzero}$ be
the set of squares of the unit group $(\Z/\Mzero\Z)^\times$.
\begin{enumerate}[label=(\roman*)]
\item $\Qsqmod{\Mzero}=\{r\in\Z/\Mzero\Z:\ r\equiv1\pmod8,\ r\bmod p\in\Qsq\
 \text{for every prime }p\mid d_0\}$. Every element of $\Qsqmod{\Mzero}$ is a
 strict square class modulo $\Mzero$ (\cref{def:strict-square-class}).
\item Under the Chinese remainder isomorphism
 $\Z/\Mzero\Z\cong\Z/8\Z\times\prod_{p\mid d_0}\Fp$, the uniform probability on
 $\Qsqmod{\Mzero}$ is the product of the point mass at $1$ in $\Z/8\Z$ and the
 uniform probabilities on the sets $\Qsq$, $p\mid d_0$. In particular
 $|\Qsqmod{\Mzero}|=\prod_{p\mid d_0}(p-1)/2$, and if $3\mid d_0$ the factor at
 $3$ is the point mass at $1$.
\end{enumerate}
\end{lemma}

\begin{proof}
Since $d_0$ is odd and squarefree, the Chinese remainder theorem gives
$\Z/\Mzero\Z\cong\Z/8\Z\times\prod_{p\mid d_0}\Fp$ and
$(\Z/\Mzero\Z)^\times\cong(\Z/8\Z)^\times\times\prod_{p\mid d_0}\Fp^\times$.
The squares of a product of groups are the products of the squares. Every
odd square is $\equiv1\pmod8$, so the squares of $(\Z/8\Z)^\times=\{1,3,5,7\}$
form $\{1\}$, and the squares of $\Fp^\times$ form $\Qsq$; this is the
description in (i), and its elements are $\equiv1\pmod8$ and nonzero squares
modulo every prime of $d_0$, which is the definition of a strict square
class. The uniform probability on a product set is the product of the
uniform probabilities, $|\Qsq|=(p-1)/2$, and $\Qsq[3]=\{1\}$; this is (ii).
\end{proof}

\begin{definition}[The coupled index]\label{def:hyb:coupled}
Let $\mathcal E_k$ be the query family of (CP0), $\Lsc\ge1$, $N=\Kint\Lsc$,
$A\subseteq[N]$ square-difference-free, and let $I_b$, $b\in\Z/\Kint\Z$, be
the intervals of \cref{def:record:functional}: consecutive intervals of
length $\Lsc$ partitioning $[N]$, each element of $I_b$ smaller than each
element of $I_{b+1}$ for $1\le b<\Kint$, while $b+1$ wraps $I_{\Kint}$ to
$I_1$.
\begin{enumerate}[label=(\alph*)]
\item\emph{Index space.} $\Omega=\prod_{e\in\mathcal E_k}\Omega_e$ with
 $\Omega_e=\Z/\Kint\Z\times\Z/\Mzero\Z$ and the uniform probability measure;
 $\omega=(\omega_e)_e$, $\omega_e=(b_e,s_e)$. By the Chinese remainder theorem
 $\Omega=\prod_{j\in\mathcal K}\Omega_j$ is a product of cyclic groups indexed
 by the \emph{coordinates} $\mathcal K=\bigsqcup_e\mathcal K_e$: the interval
 coordinate $b_e$ and the prime-power residues
 $s_e\bmod p^{h_p}$ for $p\mid\Mzero$.
\item\emph{Half-functions.} For $\omega\in\Omega$ and $x\in G_N^{V_+}$,
 \begin{align*}
  X_\omega(x)&=\prod_{e}g^{I_{b_e},\Mzero,s_e}(x_{u_e})
  \prod_{v\in V_+\setminus\{u_e\}_e}F_A(x_v),\\
  Y_\omega(x)&=\prod_{e}g^{I_{b_e},\Mzero,s_e}(x_{\iota v_e})
  \prod_{v\in V_+\setminus\{\iota v_e\}_e}F_A(x_v),
 \end{align*}
 and $\bar X(x)=\prod_{v\in V_+}F_A(x_v)$. (The half-function $Y_\omega$ always
 carries an index; it is not one of the energies $Y_\kappa$.)
\item\emph{Coupling.} Independently for every query and prime, choose
 the specified square increment $r_{e,p}\bmod p^{h_p}$; combine them by
 the Chinese remainder theorem to obtain $r_e$. The coupled index is
 $\omega'=(b+1,s+r)$. Its Markov operator
 $T_ef(b_e,s_e)=\E_{r_e}f(b_e+1,s_e+r_e)$ is the tensor product of the
 interval shift and the coordinate convolutions. Let $P_e$ average over
 the entire index $\omega_e$, interval coordinate included. All the
 coordinate transitions are doubly stochastic. The finite laws of the
 square increments $r_{e,p}$ used in the application are given in
 \cref{lem:imp:padic-envelope}; the uniform
 unit-square law of \cref{lem:func:unitsq} is a special alternative.
\item\emph{Coupled and signed forms.} For a finite signed measure $\kappa$ on
 $G_N^V$ and $\mathcal S\subseteq\mathcal E_k$ put
 \[
  I_{\kappa,\mathcal S}=\E_\omega\,Q_\kappa\Bigl(X_\omega,
  \Bigl(\bigotimes_{e\in\mathcal S}T_e\bigotimes_{e\notin\mathcal S}P_e\,Y
  \Bigr)_\omega\Bigr),\qquad
  J_\kappa=\E_\omega\,Q_\kappa\bigl(X_\omega,(\mathcal TY)_\omega\bigr),\qquad
  \mathcal T=\bigotimes_e(2P_e-T_e),
 \]
 the operators acting on the index variable of $Y_\omega$; they commute with
 the field labels. $J_\kappa$ is the \emph{signed index form}.
\item\emph{Depth law.} On a residue coordinate modulo $p^{h_p}$ let
 $E_{p,j}$ be conditional expectation given the residue modulo $p^j$,
 $0\le j\le h_p$, with $E_{p,0}$ the average. Choose decreasing tails
 $1=\delta_{p,0}\ge\delta_{p,1}\ge\cdots\ge\delta_{p,h_p}\ge0$
 bounding the transition norm on the exact-conductor spaces
 $\operatorname{ran}(E_{p,j}-E_{p,j-1})$. Set
 $\delta_{p,h_p+1}=0$ and choose a random depth with
 $\Pr(J_{e,p}\ge j)=\delta_{p,j}$, independently across primes and
 queries. Always reveal the interval. Write $\Jc$ for this whole
 depth vector and $Z^{\Jc}$ for the corresponding conditional
 expectation. Bernoulli inclusion is the special case of a chain
 containing only the constant and full-coordinate projections.
\item\emph{Energies.} For a nonnegative measure $M$ on $G_N^V$ with
 $M\succeq_\iota0$ the \emph{single-half random-inclusion energies} of $X$ and
 $Y$ under $M$ are
 \[
  V_M(X)=\E_{\Jc}\E_\omega\,Q_M\bigl(X^{\Jc}_\omega,X^{\Jc}_\omega\bigr),\qquad
  V_M(Y)=\E_{\Jc}\E_\omega\,Q_M\bigl(Y^{\Jc}_\omega,Y^{\Jc}_\omega\bigr),
 \]
 the two summands of the energy of the one-cut coupling inequality
 (\cref{lem:hyb:coupling}) without its factor $\frac12$.
\item\emph{Conductor moment.} Put
 $D_{\Jc_e}=\prod_{p\mid\Mzero}p^{J_{e,p}}$. With $\sq(D)$ the least
 square divisible by $D$, the one-query moment is
 \[
  R_I=\E_{\Jc}\bigl(\Kint\,\sq(D_{\Jc_e})\bigr)^{2a/d}.
 \]
 It does not depend on $e$. Depths are held fixed through every fold
 and structured-energy estimate, and averaged only afterwards.
\end{enumerate}
\end{definition}

The coupled forms place the two ends of every query of $\mathcal S$ on the
residue classes $s_e$ and $s_e+r_e$ of the intervals $I_{b_e}$ and
$I_{b_e+1}$. Indeed, by the bilinear identity \eqref{eq:cc:bilinear},
$Q_\kappa(X_\omega,Y_{\omega''})=\Lambda_\kappa((k_v)_{v\in V})$ with
$k_{u_e}=g^{I_{b_e},\Mzero,s_e}$, $k_{v_e}=g^{I_{b''_e},\Mzero,s''_e}$ and
$k_v=F_A$ otherwise, because the slot $\iota(\iota v_e)=v_e$ carries the
factor of $Y_{\omega''}$ at $\iota v_e$; and an averaged index $e\notin\mathcal
S$ leaves $F_A$ at both of its ends. In $V_M(Y)$ the map $Y$ is evaluated at
a uniform index; since $\omega'$ is uniform as well (the $T_j$ are doubly
stochastic), this is the natural normalization. The depth law
gives the exact norm envelope on each conductor vector
$\mathbf j$: the envelope $\prod_p\delta_{p,j_p}$ is exactly the
probability that all the depths $j_p$ are revealed.
The index-support lemma, with its filtration form below, applies to
both the transition and its adjoint.
By \cref{lem:twointerval}, $X^{\Jc}_\omega$
carries $g^{I_{b_e},D_{\Jc_e},s_e\bmod D_{\Jc_e}}$ at $u_e$ and $F_A$ at the
other slots of $V_+$, and likewise $Y^{\Jc}_\omega$ at the slots $\iota v_e$
(\cref{cor:index-support-frame}(ii)).

\begin{lemma}[Conductor moment]\label{lem:hyb:moment}
For $w\ge0$ and every query $e$,
\[
 \E_{\Jc}(\Kint\,\sq(D_{\Jc_e}))^w
 =\Kint^w\prod_{p\mid\Mzero}
 \left[1+(p^{2w}-1)\sum_{l=1}^{\lceil h_p/2\rceil}
 p^{2w(l-1)}\delta_{p,2l-1}\right].
\]
The expectation of the product over $k$ queries is the $k$-th power
of this expression. Every factor is nondecreasing in $w$ and in the
chosen tails. For the earlier squarefree law $\Mzero=8d_0$, with
$8$ always revealed and Bernoulli odd-prime probabilities $\delta_p$,
the expression is $(16\Kint)^w\prod_{p\mid d_0}(1+\delta_p(p^{2w}-1))$.
\end{lemma}
\begin{proof}
The local least square is $p^{2\lceil J_{e,p}/2\rceil}$.
For an integer-valued nonnegative variable $J$,
\[
 \E z^{\lceil J/2\rceil}
 =1+(z-1)\sum_{l\ge1}z^{l-1}\Pr(J\ge2l-1).
\]
Apply this finite summation-by-parts identity with $z=p^{2w}$;
independence across primes and queries proves the assertions.
The Bernoulli specialization follows from $\sq(8)=16$ and
$\sq(p)=p^2$. In the selected construction $w=2a/72=a/36$
and the finite-depth moment is evaluated in \cref{prop:imp:padic-price}.
\end{proof}

\paragraph{The signed index form.}
Expanding the product of the commuting operators $2P_e-T_e$,
$\bigotimes_e(2P_e-T_e)=\sum_{\mathcal S}(-1)^{|\mathcal S|}2^{k-|\mathcal
S|}\bigotimes_{e\in\mathcal S}T_e\bigotimes_{e\notin\mathcal S}P_e$, so for
every finite signed measure $\kappa$
\begin{equation}\label{eq:record:expansion}
 J_\kappa=\sum_{\mathcal S\subseteq\mathcal E_k}
 (-1)^{|\mathcal S|}2^{k-|\mathcal S|}I_{\kappa,\mathcal S},
 \qquad I_{\kappa,\varnothing}=Y_\kappa .
\end{equation}
The second identity holds because every unused index averages its slot
function to the full lift: $\E_{b,s}g^{I_b,\Mzero,s}=F_A$ for uniform $(b,s)$
(\cref{lem:twointerval}), and $X_\omega$, $Y_\omega$ are products of functions
of the independent $\omega_e$ at the distinct slots $u_e$, respectively $\iota
v_e$; so $(\bigotimes_eP_eY)_\omega=\E Y=\bar X$, $\E X=\bar X$, and
$I_{\kappa,\varnothing}=Q_\kappa(\E X,\bar X)=Y_\kappa$ by bilinearity. The constant coefficient $2^k$ of
$\mathcal T$ is the factor $2^k$ of the closing: the empty-support component
of $J_\nu$ is $Y_\nu$ itself.

\begin{remark}[The choice of the coupling]\label{rem:hyb:optimal}
Nothing in the proofs depends on the following two facts. No doubly
stochastic kernel on $\F_p$ supported on nonzero square shifts has a smaller
centred norm than the uniform square shift (\cref{lem:square-shift}(b)). Reusing one index for
two queries generally correlates their inputs, and the product mean identity
$\E X=\bar X$ then has no justification; this is why the queries have
pairwise disjoint endpoints and independent indices.
\end{remark}

\subsection{Hypotheses (CP1)--(CP8) on the three laws}\label{subsec:cp:hypotheses}

The hypotheses are conditions, atom by atom, on the reference law
$\mu$, the first-stage law $\lambda$ and the signed comparison law
$\nu$. We first fix the words \emph{record}, \emph{diagonal} and
\emph{bracket}, then state (CP1)--(CP8) in order.

\begin{definition}[Records, diagonals and brackets]\label{def:cp:records}
A \emph{record} of an atom $\Blk$ is a nonnegative measure on $X_{\Blk}^V$,
and a \emph{positive record} is a record whose normalization is admissible
(\cref{def:admissible}); in particular its one-slot marginals are uniform. A
record is \emph{diagonal at a cut} $c$ if it is fixed samplewise by the
reflection of $c$; at $\iota$ we say \emph{reserved diagonal}. A diagonal
record has a nonnegative diagonal cut matrix at its cut, so it is
$\succeq_c0$ (\cref{lem:commoncut}(i)). For a probability law $m$ on
$X_{\Blk}^{V_+}$, $\operatorname{Diag}_\iota(m)$ is its \emph{diagonal copy}
(\cref{lem:commoncut}(i)): the plus half has law $m$ and the minus half is its
reflection; its cut matrix is $\operatorname{diag}(m)$. For a measure
$\zeta$ with image $\iota_*\zeta$ under the reflection,
$\frac12(\zeta+\iota_*\zeta)$ and $\frac12(\zeta-\iota_*\zeta)$ are the
\emph{symmetric} and \emph{skew parts} of $\zeta$ at $\iota$; their cut
matrices are the symmetric and the antisymmetric part of $C_\iota(\zeta)$
(\cref{rem:cut-symmetries}(a)). A \emph{bracket} at $\iota$ is a bracket of
measures
$\bigl(\begin{smallmatrix}M&K\\ K^*&M'\end{smallmatrix}\bigr)\succeq_\iota0$
in the sense of \cref{subsec:pos:brackets}, a positive semidefinite block
matrix of cut matrices at $\iota$. It gives
$|Q_K(x,y)|\le Q_M(x,x)^{1/2}Q_{M'}(y,y)^{1/2}$, brackets tensorize over
atoms, and for $K$ symmetric at $\iota$,
$\bigl(\begin{smallmatrix}M&K\\ K&M\end{smallmatrix}\bigr)\succeq_\iota0$ if
and only if $M\pm K\succeq_\iota0$ (\cref{lem:record:block-gram}(a),(b),(d)).
\end{definition}

\begin{definition}[(CP1) Reference]\label{def:cp1}
The \emph{reference} at scale $N$ is a product
$\mu=\bigotimes_{\Blk\in\mathcal A_N}\mu_{\Blk}$ of probability measures that
satisfies \ref{it:record:R1}--\ref{it:record:R3} of
\cref{hyp:record:reference}, with \ref{it:record:R2} at every cut of
$\mathfrak F$, and the part of \ref{it:record:R4} that concerns $\mu$: at
every tail prime $p\ge\Pfive$, $\mu_p$ is $(\gamma,v)$-mixing for every
slot $v$.
\end{definition}

Reflection positivity is required at every cut of the family, not only along
the sorting sequence (see the remark after \cref{hyp:record:reference}): the
first comparison and the bracket energies of the next section, and the
restorations (\cref{lem:record:restoration}) and companion energies
(\cref{lem:record:companion-energy}) of Part~V, use chamber cuts outside the
sorting sequence.

\begin{definition}[(CP2) First stage]\label{def:cp2}
A set $\mathcal A^{\rm fs}\subseteq\mathcal A$ of \emph{first-stage atoms},
a finite set $\mathfrak T$ of \emph{tags} with a cut $c_j\in\mathfrak F$ for
every $j\in\mathfrak T$, and, for $\Blk\in\mathcal A^{\rm fs}$, signed
measures $\rho_{\Blk}$ and $\delta_{E,\Blk}$ on $X_{\Blk}^V$ ($E$ anticipates
the application, where the tagged differences below are those of the
matching records; see \cref{rem:cp:tags}) such that:
\begin{enumerate}[label=(\alph*)]
\item $\rho_{\Blk}-\mu_{\Blk}=\sum_ic_i\kappa_i$ is a finite real combination
 of positive records $\kappa_i$, with \emph{variation price}
 $p_{\Blk}=\sum_i|c_i|\kappa_i(1)\,q_{\Blk}^{2a}$, and
 $\rho_{\Blk}\succeq_{\vartheta_0}\mu_{\Blk}$ at one cut $\vartheta_0\in\mathfrak F$,
 the same for all first-stage atoms;
\item $\delta_{E,\Blk}=\sum_i(\kappa'_{\Blk,i}-\kappa_{\Blk,i})$ is a finite
 sum of differences of positive records of equal mass; each difference
 carries a tag $j(\Blk,i)\in\mathfrak T$ and has a sharp bound relative to
 $\kappa_{\Blk,i}$ at the cut $c_{j(\Blk,i)}$ of its tag
 (\cref{def:price}(b)); $d_{\Blk,i}$ is its sharp price and
 $v_{\Blk,i}=\bigl(\kappa_{\Blk,i}(1)+\kappa'_{\Blk,i}(1)\bigr)q_{\Blk}^{2a}$
 its \emph{variation price}, the price of both positive parts of the
 difference.
\end{enumerate}
The \emph{first stage} is $\lambda=\bigotimes_{\Blk}\lambda_{\Blk}$ with
$\lambda_{\Blk}=\rho_{\Blk}+\delta_{E,\Blk}$ on $\mathcal A^{\rm fs}$ and
$\lambda_{\Blk}=\mu_{\Blk}$ otherwise. Its totals are $P=\sum_\Blk p_{\Blk}$,
the \emph{tag totals}
\[
 d_j=\sum_{(\Blk,i):\,j(\Blk,i)=j}d_{\Blk,i},\qquad
 v_j=\sum_{(\Blk,i):\,j(\Blk,i)=j}v_{\Blk,i}\qquad(j\in\mathfrak T),
\]
$D=\sum_jd_j$ and $V=\sum_jv_j$, and its \emph{cost} is
\[
 L_0=e^P\tagser(d,v),\qquad
 \tagser(d,v)=\sum_{\boldsymbol k\in\Z_{\ge0}^{\mathfrak T}\setminus\{0\}}\
 \min_{j:\,k_j\ge1}\frac{d_j^{k_j}\prod_{j'\ne j}v_{j'}^{k_{j'}}}
 {\prod_{j'\in\mathfrak T}k_{j'}!},
\]
where $d=(d_j)_j$, $v=(v_j)_j$ and $\boldsymbol k=(k_j)_{j\in\mathfrak T}$
runs over the nonzero \emph{count vectors}.
\end{definition}

\begin{remark}[Tags]\label{rem:cp:tags}
The series $\tagser$ has nonnegative terms, and it is nondecreasing in every
argument, a minimum of nondecreasing functions being nondecreasing. It
converges: with $m_j=\max(d_j,v_j)$ every term at $\boldsymbol k$ is at most
$\prod_jm_j^{k_j}/k_j!$, so $\tagser(d,v)\le e^{\sum_jm_j}-1$. If
$d_j\le v_j$ for every tag, then $\tagser(d,v)\le De^V$, so $L_0\le
De^{P+V}$, the cost of expanding with one sharp difference per term: bound
the minimum at $\boldsymbol k$ by the term of one tag $j$ with $k_j\ge1$, use
$d_j^{k_j}\le d_jv_j^{k_j-1}$ and $k_j!\ge(k_j-1)!$, and sum over the
$\boldsymbol k$ assigned to each $j$. The condition $d_j\le v_j$ can always
be arranged when the two positive parts of every difference have the same
half marginals in every shared label: a difference of two probability
couplings with common half marginals satisfies \eqref{eq:record:sharp} with
the constant $2$, each coupling being a contraction between its marginal
spaces, so its sharp constant may be taken at most $2$ and then
$d_{\Blk,i}\le v_{\Blk,i}$. In that sense tags only sharpen. A tag collects
the differences that are sharp at a common cut, and all of them are applied
sharply together
(\cref{lem:sharp-first}); with a single tag the series is
$\tagser=e^{d_0}-1$, and its gain over $De^V$ comes from the differences being
paid sharply at once instead of one at a time. In Part~V there are two tags,
the two folded matching records of the one-stage repair, each sharp at its
own first cut, with $d_0=d_1=D/2$ and $v_0=v_1=V/2$
(\cref{lem:record:one-stage-folded-e,cor:two-cut-series}).
\end{remark}

\begin{definition}[(CP3) Gram bracket]\label{def:cp3}
The \emph{bracket} is $\Nbr=\bigotimes_{\Blk}\Nbr_{\Blk}$ with
$\Nbr_{\Blk}=\mu_{\Blk}$ off the first-stage atoms and
$\Nbr_{\Blk}=\Nbr_{0,\Blk}+\delta_{E,\Blk}$ on them, where
$\Nbr_{0,\Blk}-\mu_{\Blk}$ is a finite sum of positive reserved-diagonal
records, of total price $u_{\Blk}$, such that for every atom $\lambda_{\Blk}$
is symmetric at $\iota$, $\Nbr_{\Blk}$ is nonnegative, and
\[
 \begin{pmatrix}\Nbr_{\Blk}&\lambda_{\Blk}\\ \lambda_{\Blk}&\Nbr_{\Blk}
 \end{pmatrix}\succeq_\iota0,\qquad\text{equivalently}\qquad
 \Nbr_{\Blk}\pm\lambda_{\Blk}\succeq_\iota0 .
\]
We also write $\Nbrz=\bigotimes_\Blk\Nbr_{0,\Blk}$ for the
record-difference-free part of the bracket, with
$\Nbr_{0,\Blk}=\mu_\Blk$ off the first-stage atoms. Its totals
are $U_0=\sum_\Blk u_{\Blk}$ and
$R_{\rm err}=e^{U_0}\tagser(d,v)$, with the tags and prices of (CP2)(b), and
the \emph{energy factor} of the bracket is
\[
 B=e^{U_0}+R_{\rm err}.
\]
\end{definition}

The bracket is the substitute for reflection positivity of $\lambda$ at the
reserved cut: $\lambda$ need not be positive there, but it is dominated in
absolute value by a positive measure whose energies are priced. An energy
with outside law $\Nbr$ on the atoms outside its selected set $X$ costs at
most $B$ times its price (\cref{lem:bracket-energy}): $e^{U_0}$ for the
part without record differences and $R_{\rm err}$ for the rest.

\begin{definition}[(CP4) Comparison law]\label{def:record:comparison-law}
\label{def:cp4}
A \emph{comparison law} is a product $\nu=\bigotimes_{\Blk}\nu_{\Blk}$ of
finite signed measures on $X_{\Blk}^V$, each a finite real combination
$\nu_{\Blk}=\sum_{i\in I_\Blk}c_{\Blk,i}\nu_{\Blk,i}$ of nonnegative measures
$\nu_{\Blk,i}$, its \emph{components}, with $\nu_\Blk(1)>0$, and with
$|I_\Blk|=1$, $c_{\Blk,1}=1$ at all but finitely many atoms of $\mathcal A$,
such that:
\begin{enumerate}[label=(\alph*)]
\item\emph{Acceptance.} Every query of $\mathcal E_k$ is accepted by every
 component: $\nu_{\Blk,i}(1)>0$ and, $\nu_{\Blk,i}$-almost surely, for every
 $e$, the difference $z_{v_e}-z_{u_e}$ is a nonzero square at every prime of
 $\Blk$ with $e_p=1$, and at a prefix prime it is nonzero, of even valuation
 $2j<e_p$, with a square unit part. All one-slot marginals of every component
 are multiples of the uniform law. With the normalized components
 $\bar\nu_{\Blk,i}=\nu_{\Blk,i}/\nu_{\Blk,i}(1)$ as the accepted parts,
 \ref{it:record:R5} and the part of \ref{it:record:R4} that concerns
 $\bar\nu$ hold, the latter for the queries of $\mathcal E_k$ (in Part~V,
 $\mathcal E_k=\mathcal E_6$).
\item\emph{Records and brackets.} We have
 \[
  \nu_{\Blk}-\lambda_{\Blk}=\sum_{t\in\mathcal T_{\Blk}}
  \bigl(\zeta^{\rm s}_{\Blk,t}+\zeta^{\rm a}_{\Blk,t}\bigr),
 \]
 a finite sum over the records $t$ of $\Blk$ that the comparison law changes
 (repairs, replaces or deletes), where $\zeta^{\rm s}_{\Blk,t}$ and
 $\zeta^{\rm a}_{\Blk,t}$ are the symmetric and skew parts at $\iota$ of the
 change of $t$; and there are reserved-diagonal probability laws
 $D^{\rm s}_{\Blk,t}$, $D^{\rm a}_{\Blk,t}$ with uniform one-slot marginals
 (diagonal copies of actual half marginals of $t$, in the applications) and
 numbers $\tilde\pi^{\rm s}_{\Blk,t},\tilde\pi^{\rm a}_{\Blk,t}\ge0$ with
 \[
  \begin{pmatrix}\tilde\pi^{\rm s}_{\Blk,t}D^{\rm s}_{\Blk,t}&\zeta^{\rm s}_{\Blk,t}\\
  \zeta^{\rm s}_{\Blk,t}&\tilde\pi^{\rm s}_{\Blk,t}D^{\rm s}_{\Blk,t}\end{pmatrix}
  \succeq_\iota0,\qquad
  \begin{pmatrix}\tilde\pi^{\rm a}_{\Blk,t}D^{\rm a}_{\Blk,t}&\zeta^{\rm a}_{\Blk,t}\\
  \zeta^{{\rm a}*}_{\Blk,t}&\tilde\pi^{\rm a}_{\Blk,t}D^{\rm a}_{\Blk,t}\end{pmatrix}
  \succeq_\iota0 .
 \]
 The diagonal entries are the \emph{norm envelopes} of the change; their
 coefficients carry no price. The prices are
 $\pi^\epsilon_{\Blk,t}=\tilde\pi^\epsilon_{\Blk,t}q_{\Blk}^{2a}$
 ($\epsilon\in\{{\rm s},{\rm a}\}$), $s_{\Blk}=\sum_t\pi^{\rm s}_{\Blk,t}$,
 $a_{\Blk}=\sum_t\pi^{\rm a}_{\Blk,t}$, and the totals are
 \[
  S=\sum_{\Blk}s_{\Blk},\qquad A=\sum_{\Blk}a_{\Blk}.
 \]
\item\emph{Selections.} A \emph{selection} $\omega$ is a nonempty finite set
 $X\subseteq\mathcal A_N$ of atoms with a record $t_{\Blk}\in\mathcal
 T_{\Blk}$ and a type $\epsilon_{\Blk}\in\{{\rm s},{\rm a}\}$ for every
 $\Blk\in X$, the number of skew types being even. Its signed measure, its
 monomial and its term are
 \[
  \zeta_\omega=\bigotimes_{\Blk\in X}\zeta^{\epsilon_{\Blk}}_{\Blk,t_{\Blk}},\qquad
  \pi_\omega=\prod_{\Blk\in X}\pi^{\epsilon_{\Blk}}_{\Blk,t_{\Blk}},\qquad
  I_\omega=Y_{\zeta_\omega\otimes\lambda_{-X}},\qquad
  \lambda_{-X}=\bigotimes_{\Blk\notin X}\lambda_{\Blk}.
 \]
\end{enumerate}
The \emph{variation factor} of $\nu$ is
$\mathfrak v_\nu=\prod_{\Blk\in\mathcal A}|\nu|_\Blk/\nu_\Blk(1)$, with
$|\nu|_\Blk=\sum_{i\in I_\Blk}|c_{\Blk,i}|\,\nu_{\Blk,i}(1)$, and
$n_\nu=\prod_{\Blk\in\mathcal A}|I_\Blk|$ is the number of its \emph{product
components} $\bigotimes_\Blk c_{\Blk,i_\Blk}\nu_{\Blk,i_\Blk}$. Both are
finite, do not depend on the scale, and equal one when every $\nu_\Blk$ is
represented by a single nonnegative component with coefficient one;
every factor $|\nu|_\Blk/\nu_\Blk(1)$ is at least one.
\end{definition}

Summing the brackets of (b) over $t$ gives blockwise brackets with the
diagonal laws $D^\epsilon_{\Blk}=\sum_t\tilde\pi^\epsilon_{\Blk,t}
D^\epsilon_{\Blk,t}/\sum_t\tilde\pi^\epsilon_{\Blk,t}$ and the coefficients
$\sum_t\tilde\pi^\epsilon_{\Blk,t}$. A term of the expansion of
$\bigotimes_{\Blk}(\lambda_{\Blk}+\sum_t(\zeta^{\rm s}_{\Blk,t}+\zeta^{\rm
a}_{\Blk,t}))$ chooses at most one record and one type at each atom; since
$\lambda_{\Blk}$ is symmetric at $\iota$, a term with an odd number of skew
factors has a skew cut matrix and vanishes at $(\bar X,\bar X)$. This is why
selections carry an even number of skew types. In the construction of
Part~V, $\nu_{\Blk}$ is normalized by the normalizing mass of $\mu_{\Blk}$, so
that retained records keep their weights. In Part~V the middle blocks carry
signed combinations, through their fold records (three components each,
\cref{def:record:mixture}); every other atom law is nonnegative. The signs
enter the theorem only through the cancellation of the accepted queries,
which is applied to each product component separately
(\cref{lem:record:wrap-expansion}); the brackets of (b) and every energy
bound concern the record changes and the envelopes, not the components.

\begin{remark}[Charging rules]\label{rem:cp:charging}
(i) Let a record $w\kappa$, with $\kappa$ a probability measure symmetric at
$\iota$ and $m$ its half marginal, be deleted, so that $-w\kappa$ is its
change. Then $\zeta^{\rm s}=-w\kappa$, $\zeta^{\rm a}=0$, $D^{\rm
s}=\operatorname{Diag}_\iota(m)$ and $\tilde\pi^{\rm s}=w$ satisfy (b): a
deleted positive record is charged its mass in $\tilde\pi^{\rm s}$ only.
Indeed $|Q_\kappa(f,f)|\le\int|f(z_+)||f(\iota z_-)|\,d\kappa\le
\|f\|_{L^2(m)}^2=Q_{\operatorname{Diag}_\iota(m)}(f,f)$, where $\iota z_-$
denotes the minus-half labels read on $V_+$, so
$\operatorname{Diag}_\iota(m)\pm\kappa\succeq_\iota0$.
(ii) Let a record $w\Pi$ be replaced by $w\Pi'$, where $\Pi,\Pi'$ are
probability measures with the same half marginal $m$ on both halves of
$\iota$, and let $\Delta$ be the operator on $L^2(m)$ with
$\ip\varphi{\Delta\psi}_m=\int\varphi(z_+)\psi(\iota z_-)\,d(\Pi'-\Pi)(z)$, with
symmetric and skew parts $\Delta^{\rm s},\Delta^{\rm a}$. Then (b) holds with
$D^{\rm s}=D^{\rm a}=\operatorname{Diag}_\iota(m)$,
$\tilde\pi^{\rm s}=w\|\Delta^{\rm s}\|$ and $\tilde\pi^{\rm a}=w\|\Delta^{\rm
a}\|$: such a repair is charged its weight times its operator norms. Indeed
the cut matrix of $\Pi'-\Pi$ is $\operatorname{diag}(m)\Delta$ and that of
$\operatorname{Diag}_\iota(m)$ is $\operatorname{diag}(m)$, so the form of the
symmetric bracket at $(f,g)$ is
$c\|f\|_m^2+2\ip f{\Delta^{\rm s}g}_m+c\|g\|_m^2$, $c=\|\Delta^{\rm s}\|$, which
is nonnegative; likewise for the skew part.
(iii) \emph{Skew-only charging.} In (ii) let $\Pi$ be symmetric at $\iota$.
Then $\Delta^{\rm a}$ is the skew part of the operator $\Delta'$ of $\Pi'$
alone, $\ip\varphi{\Delta'\psi}_m=\int\varphi(z_+)\psi(\iota z_-)\,d\Pi'(z)$:
the operator of $\Pi'-\Pi$ is $\Delta'-\Delta_\Pi$, with $\Delta_\Pi$ the
operator of $\Pi$, which is self-adjoint. So a replacement of a symmetric
record is charged in the skew price only by the skew part of the new law,
$\tilde\pi^{\rm a}=w\|\Delta'^{\rm a}\|$, whatever the symmetric record it
replaces; and a deleted symmetric record has no skew charge, as in (i).
\end{remark}

\begin{definition}[(CP5) Class bounds]\label{def:cp5}
Optionally: for every atom $\Blk$ finitely many numbers $y_{\Blk,j}\ge0$
($j\in J_\Blk$), the \emph{auxiliary prices}, each attached to one record
$t\in\mathcal T_\Blk$ of (CP4)(b) and one type $\epsilon\in\{{\rm s},{\rm a}\}$
(in Part~V, operator prices of parts of the record changes: raw-arm, leading
and remainder prices, \cref{lem:middle-comparison-prices}); a family of
pairwise disjoint \emph{classes} $\mathcal C$ of selections, each selection
of a class having at most $n_{\max}$ atoms, with $n_{\max}$ a fixed integer
independent of the scale, of $A$ and of the class; and for every
selection $\omega$ of a class a polynomial $\phi_\omega$ with nonnegative
coefficients in the auxiliary prices, its \emph{class bound}, each of whose
monomials is a product, over the atoms $\Blk$ of $\omega$, of one auxiliary
price attached to the record and type $(t_\Blk,\epsilon_\Blk)$ of $\omega$ at
$\Blk$; such that robustly and uniformly in $\mathcal C$
\[
 \sum_{\omega\in\mathcal C}|I_\omega|\le\sum_{\omega\in\mathcal C}
 \bigl(\phi_\omega+o(1)\pi_\omega\bigr)Y .
\]
A \emph{class cost} $0\le c_{\mathcal C}\le B$ is the case
$\phi_\omega=c_{\mathcal C}\pi_\omega$, the prices
$\pi^{\epsilon_\Blk}_{\Blk,t_\Blk}$ of the selected records serving as auxiliary
prices; its \emph{gain} is $B-c_{\mathcal C}$.
\end{definition}

Every selection costs at most $B$ times its monomial
(\cref{lem:record:record-expansion}); a class cost below $B$, or class bounds
below $B$ times the monomials of their class, record a saving proved for that
class: a variance credit (\cref{sec:variance}), a bound by
positive dominating matrices (\cref{sec:envelopes}), or both. The
requirement
$c_{\mathcal C}\le B$ is no restriction: since every class satisfies the
bound with the cost $B$, a proved cost $c\ge0$ may always be replaced by
$\min(c,B)$. In particular a cost certified below an upper bound $B'$ of $B$
gives the class cost $\min(c,B)\le c$, and the closing may be evaluated with
$c$ and $B'$ (\cref{rem:cp:monotone}). A class bound in auxiliary prices is
used when it improves, in the evaluated totals, on $B$ times the full
monomials of its class; the degree-class replacement of
\cref{lem:record:record-expansion}(v) is the case used in Part~V. Since a
monomial of a class bound contains one auxiliary price of each record of its
selection, it determines the selection; this is what lets the class bounds of
all selections be added up as one power series.

\begin{definition}[(CP6) Matching envelopes]\label{def:cp6}
For $s\in\{{\rm L},{\rm R}\}$ and every atom, nonnegative measures
$M_{s,\Blk}$ such that $M_{s,\Blk}-\Nbr_{\Blk}$ is a finite nonnegative
combination of positive reserved-diagonal records, of price $T_{s,\Blk}$,
and
\[
 \begin{pmatrix}M_{{\rm L},\Blk}&\nu_{\Blk}\\ \nu_{\Blk}^*&M_{{\rm R},\Blk}
 \end{pmatrix}\succeq_\iota0
\]
for every atom. The \emph{matching exponent} is a number $T$ with
$\sum_{\Blk}T_{s,\Blk}\le T$ for both sides $s$.
\end{definition}

\begin{definition}[(CP7) Mean energies]\label{def:cp7}
A number $m$, the \emph{mean credit}, such that robustly
\[
 \tfrac12(e_{\rm L}+e_{\rm R})\ge(m-o(1))\,Y,\qquad
 e_s=Q_{\bigotimes_{\Blk}M_{s,\Blk}}(\bar X,\bar X),
\]
where $e_s$ is the \emph{mean energy} of the envelope of the side $s$.
\end{definition}

The envelopes are allowed to differ on the two sides of the reserved cut,
and each half-function is measured only under its own envelope. The mean
energies enter the closing through the common mean of the signed index
(\cref{lem:record:common-mean}): the empty-support component of $J_\nu$ is
$Y_\nu$, and that of the envelope energies is $e_s$.

\begin{definition}[Robust bounds]\label{def:robust}
Let $z\ge0$. A least normalized counterexample $(\Lsc,A)$, for
$(a,\theta,C_0)$ with any $C_0\ge1$, is
\emph{$z$-inflated} for an argument if every structured energy $\mathfrak E$
that the argument bounds at $(\Lsc,A)$ satisfies $|\mathfrak
E|\le(1+z)\|\rho\|\,\Pi\,Y$ with its coarse price $\Pi$, and the absolute
values of its \emph{remainders} (the integrals of the form bounded in
\cref{lem:record:density-can} and the sums of the form bounded in
\cref{lem:record:density-lp} that it uses) add up to at most $zY$. An
inequality $\Xi\le(c+o(1))Y$, with $\Xi$ a quantity attached to $(\Lsc,A)$ and
$c$ a number, is \emph{robust} if it has an argument such that:
\begin{enumerate}[label=(\alph*)]
\item its remainders are a number, bounded independently of $\Lsc$ and $A$,
 of integrals of the form of \cref{lem:record:density-can}, each integrated
 against a probability law of the indices and multiplied by a coefficient
 bounded independently of $\Lsc$ and $A$, and finitely many sums of the form
 of \cref{lem:record:density-lp} whose constants $c$ and $\Pi_{\rm H}$ (in the
 notation of that lemma) do not depend on $\Lsc$ and $A$;
\item there are $z_0>0$, $\Lsc_1$ and a function $c(\cdot)$ on $[0,z_0]$,
 continuous at $0$ with $c(0)=c$, such that $\Xi\le c(z)Y$ at every
 $z$-inflated least normalized counterexample with $\Lsc\ge\Lsc_1$ and $z\le
 z_0$.
\end{enumerate}
Robust lower bounds $\Xi\ge(c-o(1))Y$ are defined in the same way. A family of
bounds is robust \emph{uniformly} if one argument, one $z_0$ and $\Lsc_1$, and
one modulus of continuity at $0$ serve for all of them; in particular the
remainders of the whole family together satisfy (a).
\end{definition}

A robust bound holds with $o(1)$ in the ordinary sense. By
\cref{prop:record:density-structured}, every least normalized counterexample
with $\Lsc$ large is $z$-inflated with
$z=u_\Lsc+(C_{\rm dc}+C_{\rm R})\Lsc^{a-\sigma}$, which tends to zero uniformly
in $A$ because $a<\sigma$ and $C_0\ge1$; here $C_{\rm R}$ bounds the
remainders allowed by (a), through
\cref{lem:record:density-can,lem:record:density-lp}, as in the proof of
\cref{thm:record:lineardensity}.

\begin{definition}[(CP8) Admissibility, robustness, summability]\label{def:cp8}
\leavevmode
\begin{enumerate}[label=(\alph*)]
\item\emph{Admissibility.} Every law named in (CP2)--(CP6) (the positive
 records of (CP2)(a), the records $\kappa_{\Blk,i},\kappa'_{\Blk,i}$, the
 records of $\Nbr_{0,\Blk}-\mu_{\Blk}$, the laws $D^\epsilon_{\Blk,t}$ and the
 records of $M_{s,\Blk}-\Nbr_{\Blk}$) is admissible, and so is every selected
 law that the reductions of this part produce from them; every reduction of
 this part satisfies (O3) and (O4) of \cref{def:admissible}.
\item\emph{Robustness.} The bounds of (CP5) and (CP7) are robust
 (\cref{def:robust}), those of (CP5) uniformly in the class.
\item\emph{Summability.} With $q_{\Blk}$ here denoting the
 product of the primes of $\Blk$ with $e_p=1$ only (so the weight
 $q_{\Blk}^{1/4}$ is $1$ at a prefix atom, whose prices in
 \cref{def:price} nevertheless carry the factor $p^{e_pa}$),
 $v_\Blk=\sum_iv_{\Blk,i}$ and $d_\Blk=\sum_id_{\Blk,i}$
 as in (CP2), $s_{\Blk},a_{\Blk}$ as in (CP4) and the auxiliary prices of
 (CP5),
 \[
  \sum_{\Blk\in\mathcal A}q_{\Blk}^{1/4}\Bigl(p_{\Blk}+v_{\Blk}+d_{\Blk}
  +u_{\Blk}+s_{\Blk}+a_{\Blk}+T_{{\rm L},\Blk}+T_{{\rm R},\Blk}
  +\sum_{j\in J_\Blk}y_{\Blk,j}\Bigr)<\infty .
 \]
\end{enumerate}
\end{definition}

Part (c) makes the totals $P$, $V$, $D$ (and the tag totals), $U_0$, $S$,
$A$, $T$ and the totals of the auxiliary prices finite and independent of the
scale, and it makes the product $\Pi_{\rm H}$ of \cref{lem:record:density-lp}
finite for the coefficients that the expansions produce: these are
post-reduction prices of the selected laws, and only the tail atoms form an
infinite family. A bound is shown to be robust by the criterion of
\cref{lem:record:rec}(i): it is built from structured-energy bounds by sums,
products and monotone operations, or it is an exact relation minus such
bounds, and every strict inequality it needs holds strictly at $z=0$.

\subsection{Theorem~C: if \texorpdfstring{$\Theta<1$}{Theta<1} then \texorpdfstring{$s(N)\le C'N^{1-a/(72-\theta)}$}{s(N) <= C' N\textasciicircum (1-a/(72-theta))}}\label{subsec:cp:theorem}

\paragraph{How to read the statement.}
Theorem~C is conditional: it assumes laws satisfying (CP0)--(CP8)
and concludes the power saving from the single numerical inequality
$\Theta<1$. The symbols in $\Theta$ are: $d=72$ slots and $k$ queries
($k=6$ in Part~V); $L_0$ the total price of $\mu\to\lambda$ (CP2);
$B$ the energy factor of the law dominating $\lambda$ (CP3);
$S,A$ the total symmetric and skew prices of $\nu-\lambda$ (CP4),
with $\hat S,\hat A$ any finite upper bounds for them; $\mathcal E$
any number at least $\Psi(\hat\pi,\hat y)$, where $\Psi$ is a
\emph{complete polynomial}---a power series with nonnegative
coefficients dominating every simultaneous-change bound, defined in
\cref{rem:cp:monotone} below and constructed in
\cref{lem:record:record-expansion} (without (CP5) families,
$\Psi(\hat\pi)=B(e^{\hat S}\cosh\hat A-1)$; the simpler display
$\mathcal E=B(e^S\cosh A-1)-D$ in \cref{subsec:intro:inequality} is
the Part~V evaluation with savings $D$ subtracted); $T$ the total
price of the dominators for $\nu$ (CP6), entering through $e^T$;
$R_I$ the conductor factor of one index (\cref{def:hyb:coupled});
$m$ the mean-energy lower bound (CP7), whose trivial admissible
value is $1-L_0$ (\cref{lem:record:domination}); and
$J,W$ as displayed. Hatted quantities are certified upper bounds:
by monotonicity (\cref{rem:cp:monotone}) the theorem may be applied
with upper bounds for prices and a lower bound for $m$. Note
$0<\theta\le1$ in (CP0): the improved exponent
$\vartheta=\FourRoutesDensity>1$ of \cref{thm:main} is not obtained
by substituting $\vartheta$ for $\theta$ here, but by the separate
induction of \crefrange{sec:density:moments-prefix}{sec:four-routes-closing}.
In (CP4) the letter $\omega$ denotes a \emph{selection} (a choice of
changed records), not the random index $\omega$ of
\cref{def:hyb:coupled}; the context (prices versus intervals)
distinguishes them.

\begin{theorem}[Theorem~C. If $\Theta<1$ then $s(N)\le C'N^{1-a/(72-\theta)}$]\label{thm:comparison}
Let comparison data satisfy (CP0)--(CP8) at every scale. Let $\hat\pi\ge\pi$
be upper prices for the records of (CP4), with totals $\hat S\ge S$ and
$\hat A\ge A$ finite; let $\hat y\ge y$ be upper auxiliary prices of (CP5),
of finite total $\sum_{\Blk,j}\hat y_{\Blk,j}$; let $\Psi$ be a complete
polynomial of \cref{lem:record:record-expansion}(ii) for the classes and
class bounds of (CP5) (for instance in the completed form (iv) or the
replaced form (v) of that lemma) that converges at $(\hat\pi,\hat y)$, and
let $\mathcal E\ge\Psi(\hat\pi,\hat y)$. Put
\begin{gather*}
 J=Be^T\Bigl(\frac{R_I}2\Bigr)^k,\qquad
 W=Be^{S+A}\Bigl[\Bigl(1+\frac{R_I}{2\Kint^{1-2\theta/d}}\Bigr)^k-1\Bigr],\\
 \Theta=\bigl(1-2^{-k}\bigr)\bigl(L_0+\mathcal E\bigr)+J+W+\frac{1-m}{2^k}.
\end{gather*}
If $\Theta<1$, there is a constant $C$, depending only on the
comparison data, such that $\Yf(\Lsc,A)\le
C\Lsc^{-a}\alpha^\theta$ for every $\Lsc\ge1$ and every square-difference-free
set $A\subseteq[\Kint\Lsc]$, $\alpha=|A|/(\Kint\Lsc)$. Consequently
\[
 s(N)\le C'N^{1-a/(d-\theta)}\qquad\text{for every }N\ge1 .
\]
\end{theorem}

The proof is in \cref{sec:comparison-proof}. Without classes, the complete
polynomial of \cref{lem:record:record-expansion}(ii) is
$\Psi=B(e^{S(x)}\cosh A(x)-1)$ and $\Psi(\hat\pi)=B(e^{\hat S}\cosh\hat
A-1)$; a class replaces $B$ times its monomials by its class bounds.

\begin{remark}[Monotonicity]\label{rem:cp:monotone}
$\Theta$ is nondecreasing in $L_0$, $\mathcal E$, $B$, $S$, $A$, $T$ and
$R_I$, and nonincreasing in $m$: $1-2^{-k}>0$, $J$ is nondecreasing in $B$,
$T$ and $R_I$, and $W$ in $B$, $S+A$ and $R_I$, all of these being
nonnegative. A complete polynomial is any power series with nonnegative
coefficients that dominates, coefficient by coefficient, the series
$\sum_\omega b_\omega$ of \cref{lem:record:record-expansion}(ii), in which a
selection in no class contributes $B$ on its monomial, a selection of a class
with a class cost that cost on its monomial, and a selection of a class with
class bounds its class bound. So every power series with nonnegative
coefficients dominating these, for instance one built from certified upper
bounds for $B$, for the class costs and for the coefficients of the class
bounds, gives an admissible $\mathcal E$ when it is evaluated at upper prices
and upper auxiliary prices. Hence the theorem holds with certified upper
bounds for $L_0$, $B$, $S$, $A$, $T$, $R_I$, the prices, the auxiliary prices,
the class costs and the coefficients of the class bounds, and a certified
lower bound for $m$, in place of the quantities themselves; this is how the
closing of Part~V evaluates $\Theta$.
\end{remark}

\begin{remark}[Reading the closing number]\label{rem:cp:reading}
The sum $L_0+\mathcal E$ bounds $|Y_\nu-Y_\mu|/Y$ up to $o(1)$, and
since the empty term of the signed-index expansion is $Y_\nu$
itself, the bound $L_0+\mathcal E$ enters $\Theta$ with the
coefficient $(2^k-1)/2^k$ (see the opening of
\cref{sec:comparison-lemmas}). In $J$ each query contributes
$R_I/2$; $W$ bounds the terms in which every query to which the
coupled transition, rather than the average, is applied in that
term wraps; and
$m>1$ lowers $\Theta$ through $(1-m)/2^k$. In the baseline
application $k=6$, $d=72$, $\theta=1$, so $1-2^{-k}=63/64$ and
$\Kint^{1-2\theta/d}=\Kint^{35/36}$ with $\Kint=\PushK$, giving
$\eps=a/71$ and $\Theta<\PushTheta$ (\cref{sec:closing}). These are
baseline values, quoted here only as an illustration; Theorem~C
itself uses no Part~V number.
\end{remark}

\paragraph{Consequences and what to read next.}
From Part~II: the slot array and its cuts (\cref{lem:slot-array},
\cref{def:sorting}), the cut calculus (\cref{def:cc:order},
\cref{eq:cc:bilinear}, \cref{rem:cut-symmetries}, \cref{lem:commoncut},
\cref{lem:record:block-gram}), the
functional and its inputs (\cref{def:record:functional},
\cref{lem:twointerval}), the reference hypotheses
(\cref{hyp:record:reference}), the square shift (\cref{lem:square-shift}), the
sector lemma (\cref{lem:sector}), and the analytic notions of
\cref{def:record:analytic-stage,def:record:counterexample,def:record:structured-energy,def:admissible},
\cref{eq:record:sharp}, \cref{lem:record:density-can,lem:record:density-lp}.
